\documentclass[10pt,a4paper]{amsart}
\usepackage{amsmath,amssymb,amsthm,graphicx}
% For commutative diagrams we use Xy-pic
\usepackage[all]{xy}

% We use 2cell for 2-commutative diagrams.
\xyoption{2cell}
% Theorem environments.
%
\theoremstyle{plain}
\newtheorem{theorem}[subsection]{Theorem}
\newtheorem{proposition}[subsection]{Proposition}
\newtheorem{lemma}[subsection]{Lemma}

\theoremstyle{definition}
\newtheorem{definition}[subsection]{Definition}
\newtheorem{example}[subsection]{Example}
\newtheorem{exercise}[subsection]{Exercise}
\newtheorem{situation}[subsection]{Situation}

\theoremstyle{remark}
\newtheorem{remark}[subsection]{Remark}
\newtheorem{remarks}[subsection]{Remarks}

\numberwithin{equation}{subsection}

% Macros
%
\def\lim{\mathop{\mathrm{lim}}\nolimits}
\def\colim{\mathop{\mathrm{colim}}\nolimits}
\def\Spec{\mathop{\mathrm{Spec}}}
\def\Hom{\mathop{\mathrm{Hom}}\nolimits}
\def\Ext{\mathop{\mathrm{Ext}}\nolimits}
\def\SheafHom{\mathop{\mathcal{H}\!\mathit{om}}\nolimits}
\def\SheafExt{\mathop{\mathcal{E}\!\mathit{xt}}\nolimits}
\def\Sch{\mathit{Sch}}
\def\Mor{\mathop{\mathrm{Mor}}\nolimits}
\def\Ob{\mathop{\mathrm{Ob}}\nolimits}
\def\Sh{\mathop{\mathit{Sh}}\nolimits}
\def\NL{\mathop{N\!L}\nolimits}
\def\CH{\mathop{\mathrm{CH}}\nolimits}
\def\proetale{{pro\text{-}\acute{e}tale}}
\def\etale{{\acute{e}tale}}
\def\QCoh{\mathit{QCoh}}
\def\Ker{\mathop{\mathrm{Ker}}}
\def\Im{\mathop{\mathrm{Im}}}
\def\Coker{\mathop{\mathrm{Coker}}}
\def\Coim{\mathop{\mathrm{Coim}}}

% Boxtimes
%
\DeclareMathSymbol{\boxtimes}{\mathbin}{AMSa}{"02}

%
% Macros for moduli stacks/spaces
%
\def\QCohstack{\mathcal{QC}\!\mathit{oh}}
\def\Cohstack{\mathcal{C}\!\mathit{oh}}
\def\Spacesstack{\mathcal{S}\!\mathit{paces}}
\def\Quotfunctor{\mathrm{Quot}}
\def\Hilbfunctor{\mathrm{Hilb}}
\def\Curvesstack{\mathcal{C}\!\mathit{urves}}
\def\Polarizedstack{\mathcal{P}\!\mathit{olarized}}
\def\Complexesstack{\mathcal{C}\!\mathit{omplexes}}
% \Pic is the operator that assigns to X its picard group, usage \Pic(X)
% \Picardstack_{X/B} denotes the Picard stack of X over B
% \Picardfunctor_{X/B} denotes the Picard functor of X over B
\def\Pic{\mathop{\mathrm{Pic}}\nolimits}
\def\Picardstack{\mathcal{P}\!\mathit{ic}}
\def\Picardfunctor{\mathrm{Pic}}
\def\Deformationcategory{\mathcal{D}\!\mathit{ef}}
\usepackage{cleveref}
\crefname{originallemma}{original lemma}{original lemmas}
\Crefname{originallemma}{Original lemma}{Original lemmas}
\makeatletter
\newlabel{algebra-definition-distance}{{02MG}{source}{}{sourcepointer-1}{}}
\newlabel{algebra-definition-distance@cref}{{[originallemma][0][]02MG}{[1][0][]source}{}{}{}}
\makeatother
\crefname{originallemma}{original lemma}{original lemmas}
\Crefname{originallemma}{Original lemma}{Original lemmas}
\makeatletter
\newlabel{algebra-lemma-order-vanishing-determinant}{{02MI}{source}{}{sourcepointer-2}{}}
\newlabel{algebra-lemma-order-vanishing-determinant@cref}{{[originallemma][0][]02MI}{[1][0][]source}{}{}{}}
\makeatother
\crefname{originallemma}{original lemma}{original lemmas}
\Crefname{originallemma}{Original lemma}{Original lemmas}
\makeatletter
\newlabel{algebra-lemma-properties-distance-function}{{02MH}{source}{}{sourcepointer-3}{}}
\newlabel{algebra-lemma-properties-distance-function@cref}{{[originallemma][0][]02MH}{[1][0][]source}{}{}{}}
\makeatother
\crefname{originallemma}{original lemma}{original lemmas}
\Crefname{originallemma}{Original lemma}{Original lemmas}
\makeatletter
\newlabel{definition-rational-equivalence}{{02RW}{source}{}{sourcepointer-4}{}}
\newlabel{definition-rational-equivalence@cref}{{[originallemma][0][]02RW}{[1][0][]source}{}{}{}}
\makeatother
\crefname{originallemma}{original lemma}{original lemmas}
\Crefname{originallemma}{Original lemma}{Original lemmas}
\makeatletter
\newlabel{divisors-definition-regular-meromorphic-ideal-denominators}{{02P1}{source}{}{sourcepointer-5}{}}
\newlabel{divisors-definition-regular-meromorphic-ideal-denominators@cref}{{[originallemma][0][]02P1}{[1][0][]source}{}{}{}}
\makeatother
\crefname{originallemma}{original lemma}{original lemmas}
\Crefname{originallemma}{Original lemma}{Original lemmas}
\makeatletter
\newlabel{divisors-lemma-regular-meromorphic-ideal-denominators}{{02P0}{source}{}{sourcepointer-6}{}}
\newlabel{divisors-lemma-regular-meromorphic-ideal-denominators@cref}{{[originallemma][0][]02P0}{[1][0][]source}{}{}{}}
\makeatother
\crefname{originallemma}{original lemma}{original lemmas}
\Crefname{originallemma}{Original lemma}{Original lemmas}
\makeatletter
\newlabel{lemma-cycles-k-group}{{02S9}{source}{}{sourcepointer-7}{}}
\newlabel{lemma-cycles-k-group@cref}{{[originallemma][0][]02S9}{[1][0][]source}{}{}{}}
\makeatother
\crefname{originallemma}{original lemma}{original lemmas}
\Crefname{originallemma}{Original lemma}{Original lemmas}
\makeatletter
\newlabel{lemma-finite-cycles-k-group}{{0FDR}{source}{}{sourcepointer-8}{}}
\newlabel{lemma-finite-cycles-k-group@cref}{{[originallemma][0][]0FDR}{[1][0][]source}{}{}{}}
\makeatother
\crefname{originallemma}{original lemma}{original lemmas}
\Crefname{originallemma}{Original lemma}{Original lemmas}
\makeatletter
\newlabel{situation-setup}{{02QL}{source}{}{sourcepointer-9}{}}
\newlabel{situation-setup@cref}{{[originallemma][0][]02QL}{[1][0][]source}{}{}{}}
\makeatother
\title{A canonical map from the Chow group of k-cycles to K\_0 of the quotient of coherent sheaves in adjacent dimensions}
\author{The Stacks Project Authors}
\date{stacks-project:chow.tex:local-snapshot-20261012; GNU FDL 1.2 or later}
\begin{document}
\maketitle
\noindent\textit{Editorial note.} All delivered mathematical text is the author's own, copied from the pinned Stacks Project chapter file chow.tex; no statement or proof step was rewritten, reordered or removed.\par\medskip
\noindent\textit{Editorial note.} Preamble provenance: the pinned chapter file opens with \textbackslash{}input\{preamble\} and carries no preamble text of its own, so the project's single Stacks-wide preamble staged at sources/stacks-chow/preamble.tex is excerpted verbatim as two spans, lines 14-18 (Xy-pic setup) and lines 154-217 (the native theorem declarations and mathematical macros). Nothing was added, removed or redefined; the wrapper class and the amsmath/amssymb/amsthm/graphicx base packages are supplied by the assembler.\par\medskip
\noindent\textit{Editorial note.} Same-chapter pointers are located at their \textbackslash{}label lines inside the pinned chapter file; cross-chapter pointers are located in the project's own tag list staged in the snapshot (sources/stacks-chow/tags.txt), at the verbatim lines carrying their canonical tags.\par\medskip
\noindent\textit{Allowed standard prerequisites (editorial).}\begin{itemize}
\item Isomorphism between k-cycles and K\_0 of the truncated coherent category. Same chapter chow.tex, label lemma-cycles-k-group, invoked at proof line 3944. Provides the map from Z\_k(X) to K\_0(Coh\_\{<=k\}(X)/Coh\_\{<=k-1\}(X)) whose composite with the quotient map is the asserted canonical map; the descent through rational equivalence is what the counted proof supplies.
\item Pushforward on K\_0 for a finite morphism of schemes. Same chapter chow.tex, label lemma-finite-cycles-k-group, invoked at proof line 3959. Lets the disjoint union W of the integral subschemes W\_j be used in place of X, so that one may assume X integral of dimension k + 1.
\item The ideal of denominators of a rational function. Divisors, labels divisors-definition-regular-meromorphic-ideal-denominators (proof line 3971) and divisors-lemma-regular-meromorphic-ideal-denominators (proof line 3982). Produces the ideal I with the two short exact sequences whose difference of classes represents div(f), which is the construction at the heart of the proof.
\item The distance function and the order of vanishing of a determinant. Algebra, labels algebra-definition-distance (proof line 4008), algebra-lemma-properties-distance-function and algebra-lemma-order-vanishing-determinant (proof lines 4014-4015); same chapter chow.tex, label definition-rational-equivalence (proof line 3952). Converts length\_A(A/I) - length\_A(A/fI) into ord\_A(f) so that the claimed cycle identity matches the coefficient of div(f).
\end{itemize}\medskip
\ifdefined\hypertarget\hypertarget{sourcepointer-1}{}\fi
\noindent\textit{Source reference.} Original lemma 02MG: sources/stacks-chow/tags.txt, source lines 3443--3443. Navigation pointer for \textbackslash{}ref\{algebra-definition-distance\} used inside the counted span (line 4008); the argument that consumes it is retained in the span.\par\medskip
\ifdefined\hypertarget\hypertarget{sourcepointer-2}{}\fi
\noindent\textit{Source reference.} Original lemma 02MI: sources/stacks-chow/tags.txt, source lines 3445--3445. Navigation pointer for \textbackslash{}ref\{algebra-lemma-order-vanishing-determinant\} used inside the counted span (line 4015); the argument that consumes it is retained in the span.\par\medskip
\ifdefined\hypertarget\hypertarget{sourcepointer-3}{}\fi
\noindent\textit{Source reference.} Original lemma 02MH: sources/stacks-chow/tags.txt, source lines 3444--3444. Navigation pointer for \textbackslash{}ref\{algebra-lemma-properties-distance-function\} used inside the counted span (line 4014); the argument that consumes it is retained in the span.\par\medskip
\ifdefined\hypertarget\hypertarget{sourcepointer-4}{}\fi
\noindent\textit{Source reference.} Original lemma 02RW: sources/stacks-chow/pinned-original-source.tex, source lines 2882--2882. Navigation pointer for \textbackslash{}ref\{definition-rational-equivalence\} used inside the counted span (line 3952); the argument that consumes it is retained in the span.\par\medskip
\ifdefined\hypertarget\hypertarget{sourcepointer-5}{}\fi
\noindent\textit{Source reference.} Original lemma 02P1: sources/stacks-chow/tags.txt, source lines 3536--3536. Navigation pointer for \textbackslash{}ref\{divisors-definition-regular-meromorphic-ideal-denominators\} used inside the counted span (line 3971); the argument that consumes it is retained in the span.\par\medskip
\ifdefined\hypertarget\hypertarget{sourcepointer-6}{}\fi
\noindent\textit{Source reference.} Original lemma 02P0: sources/stacks-chow/tags.txt, source lines 3535--3535. Navigation pointer for \textbackslash{}ref\{divisors-lemma-regular-meromorphic-ideal-denominators\} used inside the counted span (line 3982); the argument that consumes it is retained in the span.\par\medskip
\ifdefined\hypertarget\hypertarget{sourcepointer-7}{}\fi
\noindent\textit{Source reference.} Original lemma 02S9: sources/stacks-chow/pinned-original-source.tex, source lines 3757--3757. Navigation pointer for \textbackslash{}ref\{lemma-cycles-k-group\} used inside the counted span (line 3944); the argument that consumes it is retained in the span.\par\medskip
\ifdefined\hypertarget\hypertarget{sourcepointer-8}{}\fi
\noindent\textit{Source reference.} Original lemma 0FDR: sources/stacks-chow/pinned-original-source.tex, source lines 3898--3898. Navigation pointer for \textbackslash{}ref\{lemma-finite-cycles-k-group\} used inside the counted span (line 3959); the argument that consumes it is retained in the span.\par\medskip
\ifdefined\hypertarget\hypertarget{sourcepointer-9}{}\fi
\noindent\textit{Source reference.} Original lemma 02QL: sources/stacks-chow/pinned-original-source.tex, source lines 1372--1372. Navigation pointer for \textbackslash{}ref\{situation-setup\} used inside the counted span (line 3936); the argument that consumes it is retained in the span.\par\medskip
\begin{lemma}
\label{lemma-from-chow-to-K}
Let $X$ be a scheme locally of finite type over $(S, \delta)$
as in Situation \ref{situation-setup}. There is a canonical map
$$
\CH_k(X)
\longrightarrow
K_0(\textit{Coh}_{\leq k + 1}(X)/\textit{Coh}_{\leq k - 1}(X))
$$
induced by the map
$Z_k(X) \to K_0(\textit{Coh}_{\leq k}(X)/\textit{Coh}_{\leq k - 1}(X))$
from Lemma \ref{lemma-cycles-k-group}.
\end{lemma}

\begin{proof}
We have to show that an element $\alpha$ of $Z_k(X)$ which is rationally
equivalent to zero, is mapped to zero in
$K_0(\textit{Coh}_{\leq k + 1}(X)/\textit{Coh}_{\leq k - 1}(X))$.
Write $\alpha = \sum (i_j)_*\text{div}(f_j)$ as in
Definition \ref{definition-rational-equivalence}.
Observe that
$$
\pi = \coprod i_j : W = \coprod W_j \longrightarrow X
$$
is a finite morphism as each $i_j : W_j \to X$ is a closed immersion
and the family of $W_j$ is locally finite in $X$. Hence we may use
Lemma \ref{lemma-finite-cycles-k-group} to reduce to the case of $W$.
Since $W$ is a disjoint union of integral scheme, we reduce
to the case discussed in the next paragraph.

\medskip\noindent
Assume $X$ is integral of $\delta$-dimension $k + 1$.
Let $f$ be a nonzero rational function on $X$.
Let $\alpha = \text{div}(f)$. We have to show that
$\alpha$ is mapped to zero in
$K_0(\textit{Coh}_{\leq k + 1}(X)/\textit{Coh}_{\leq k - 1}(X))$.
Let $\mathcal{I} \subset \mathcal{O}_X$ be the ideal of denominators
of $f$, see Divisors, Definition
\ref{divisors-definition-regular-meromorphic-ideal-denominators}.
Then we have short exact sequences
$$
0 \to \mathcal{I} \to \mathcal{O}_X \to \mathcal{O}_X/\mathcal{I} \to 0
$$
and
$$
0 \to \mathcal{I} \xrightarrow{f} \mathcal{O}_X \to
\mathcal{O}_X/f\mathcal{I} \to 0
$$
See Divisors, Lemma
\ref{divisors-lemma-regular-meromorphic-ideal-denominators}.
We claim that
$$
[\mathcal{O}_X/\mathcal{I}]_k - [\mathcal{O}_X/f\mathcal{I}]_k =
\text{div}(f)
$$
The claim implies the element $\alpha = \text{div}(f)$ is represented by
$[\mathcal{O}_X/\mathcal{I}] - [\mathcal{O}_X/f\mathcal{I}]$
in $K_0(\textit{Coh}_{\leq k}(X)/\textit{Coh}_{\leq k - 1}(X))$.
Then the short exact sequences show that this element maps to
zero in $K_0(\textit{Coh}_{\leq k + 1}(X)/\textit{Coh}_{\leq k - 1}(X))$.

\medskip\noindent
To prove the claim, let $Z \subset X$ be an integral closed subscheme
of $\delta$-dimension $k$ and let $\xi \in Z$ be its generic point.
Then $I = \mathcal{I}_\xi \subset A = \mathcal{O}_{X, \xi}$
is an ideal such that $fI \subset A$. Now the coefficient of
$[Z]$ in $\text{div}(f)$ is $\text{ord}_A(f)$. (Of course as usual
we identify the function field of $X$ with the fraction field of $A$.)
On the other hand, the coefficient of $[Z]$ in
$[\mathcal{O}_X/\mathcal{I}] - [\mathcal{O}_X/f\mathcal{I}]$
is
$$
\text{length}_A(A/I) - \text{length}_A(A/fI)
$$
Using the distance function of
Algebra, Definition \ref{algebra-definition-distance}
we can rewrite this as
$$
d(A, I) - d(A, fI) = d(I, fI) = \text{ord}_A(f)
$$
The equalities hold by Algebra, Lemmas
\ref{algebra-lemma-properties-distance-function} and
\ref{algebra-lemma-order-vanishing-determinant}.
(Using these lemmas isn't necessary, but convenient.)
\end{proof}

\end{document}
